\section{ОСНОВНАЯ ЧАСТЬ}



\subsection{Посещение лекционных занятий}

В рамках достижения цели были прослушаны 2 лекционных занятия, в которых рассказали общие понятия проектной документации и подходы к её оформлению.

\subsubsection{Лекция 1}

На первом занятии я слушал лекцию и фиксировал ключевые определения и структуру проектной документации, 
чтобы затем опираться на них при подготовке собственных материалов.

Я понял, что проектная документация представляет собой совокупность документов, 
в которых отражаются цели и задачи проекта, сроки, риски и ресурсы.
Также я отметил, что документация нужна не только для фиксации требований, 
но и для поддержки и развития системы (например, чтобы новые участники команды могли быстрее разобраться в проекте). 
Форматы документов могут быть как электронные, так и физические.

В ходе лекции я выписал основные типы проектных документов:
\begin{enumerate}[label=\arabic*)]
  \item Техническое задание;
  \item Проектный план;
  \item График работ;
  \item Смета затрат;
  \item Отчеты о прогрессе;
  \item Документация по качеству;
  \item Документация по безопасности;
  \item Документация по разработке;
  \item Документация по тестированию;
  \item Документация по развертыванию;
  \item Документация по эксплуатации;
  \item Документация по безопасности.
\end{enumerate}

Отдельно я отметил, что документы можно классифицировать и по видам:
\begin{enumerate}[label=\arabic*)]
  \item Техническая документация;
  \item Финансовая документация;
  \item Юридическая документация;
  \item Управленческая документация;
  \item Agile-документация;
  \item Водопадная документация;
  \item Документация по DevOps;
  \item Документация по облачным сервисам.
\end{enumerate}

Дополнительно я обратил внимание на понятие конструкторской документации: 
это графические и текстовые документы, определяющие устройство изделия на каждом этапе жизненного цикла.

Далее на занятии я разобрался в структуре выпускной квалификационной работы и зафиксировал обязательные элементы:
\begin{enumerate}[label=\arabic*)]
  \item Титульный лист;
  \item Задание;
  \item Аннотация;
  \item Содержание;
  \item Список сокращений и условных обозначений (опционально);
  \item Термины и определения (опционально);
  \item Текст ВКР (введение, основная часть, заключение);
  \item Список использованных источников;
  \item Список иллюстрированного материала (опционально);
  \item Приложения (опционально).
\end{enumerate}
Я уточнил требования к содержанию введения: в нём должны быть актуальность темы, степень разработанности, решаемая проблема, цель и задачи, 
а также практическая значимость.
Кроме того, я выписал ключевые требования к оформлению (позже они были высланы отдельным документом: «ТРЕБОВАНИЯ К ВЫПУСКНЫМ КВАЛИФИКАЦИОННЫМ РАБОТАМ», код идентификации документа — ЛНАОБУЧ-СМК-03-05-2022, версия 4.0). 
Объём ВКР должен составлять 40--50 страниц текстовой части (без титульного листа, приложений и библиографии).
По списку источников я зафиксировал, что можно использовать книги, статьи, наборы данных, исходный код и техническую документацию; 
ориентир по количеству — 20--30 источников, при этом примерно \(2/3\) должны быть не старше 5 лет (на любом языке).
Также я отметил требование о размещении в приложении полного исходного кода разработанного приложения.

Из лекции я вынес перечень основных стандартов документирования:
\begin{enumerate}[label=\arabic*)]
  \item ЕСКД (ГОСТ 2) — Единая система конструкторской документации;
  \item ЕСПД (ГОСТ 19) — Единая система программной документации, регламентирует создание документации для программного обеспечения;
  \item ГОСТ 34 — стандарт на техническое задание для создания автоматизированной системы, включающей ПО, аппаратное обеспечение, людей и автоматизируемые процессы.
\end{enumerate}
Также я зафиксировал, что могут использоваться и международные стандарты: IEEE STD 830-1998, ISO/IEC/IEEE 29148, BABOK.

Отдельный блок был посвящён тому, как должна выглядеть структура технического задания. Я выписал следующие пункты:
\begin{enumerate}[label=\arabic*)]
  \item Полное и краткое название разрабатываемого ПО или проекта;
  \item Пояснение для чего, в какой области и с какой целью разрабатывается система;
  \item Документы, на основании которых производится разработка;
  \item Функций, которые выполняет разрабатываемая система или проект;
  \item Архитектура и компоненты проекта;
  \item Пользовательский интерфейсж;
  \item Требования к надёжности, безопасности и условиям эксплуатации;
  \item Какие документы разработаны;
  \item Последовательность выполнения работ;
  \item Описание процесса сдачи работы.
\end{enumerate}

Я также узнал о документах типа "программа и методики испытаний": они описывают процесс тестирования, объект и цель испытаний, а также требования к программе испытаний.

В конце я зафиксировал важное уточнение про Agile: распространено заблуждение, что документация не нужна, однако корректный подход — документации должно быть достаточно, но без избыточности.
Именно поэтому я выписал типичные артефакты Agile: backlog, user stories, критерии приёмки, definition of done, глоссарий, макеты, тестовую документацию, доски задач, роадмап и описание API.

Для себя я отметил инструменты, которые можно использовать для ведения документации: Confluence, Jira, Trello, Redmine, Notion, Miro, а также подходы, где документация создаётся кодом.

Подводя итог, лекция дала мне целостное понимание того, какие виды проектной документации существуют и какие требования предъявляются к ВКР и техническому заданию.
За счёт этого у меня сформировался ориентир по структуре и содержанию отчётных материалов, а также по стандартам и инструментам, на которые можно опираться при подготовке документации.

\subsubsection{Лекция 2}

В рамках второго занятия я подключился к платформе KTalk, где проводилась лекция и разбор отчётных материалов обучающихся.
Преподаватель пояснил, что в ходе занятия можно поднять руку и запросить разбор ошибок в отчёте конкретного человека.

Я поднял руку для получения обратной связи и получил ревью по своему отчёту: преподаватель указал места, которые необходимо исправить.
Первое замечание, которое подметил преподаватель, касалось введения: нужны были более подробные формулировки задания и этапов выполнения; во введении следует явно раскрыть содержание каждого этапа и соответствующее ему задание (для наглядности допускается оформить этапы в виде таблицы «название этапа» — «задание»).
Второе замечание касалось описания лекции 1: требовалось переписать изложение от первого лица, чтобы текст отражал мои действия на занятии (что я слушал, фиксировал и выписывал) и был связным с дальнейшими выводами в отчёте.
Следующее замечание относилось к разделу о разработке шаблона: преподаватель рекомендовал явно обосновать выбор LaTeX{} и перечислить, на какие пункты документа «Требования к выпускным квалификационным работам» я опирался при реализации требований к оформлению.
Также преподаватель отметил, что в отчёте не хватало подведения итогов по выполненным заданиям: по пункту 1.2 следовало указать, что я сформировал архив с результатами, отправил его преподавателю по электронной почте и получил отметку о выполнении задания.
Ещё одно замечание касалось того, что в пункте 1.3 необходимо явно описывать сделанную работу: 
что именно было выполнено по этапу, 
на что опирался и какой получился результат, чтобы было видно соответствие заданию практики. Так же по этому этапу мне нужно было прикрепить во вложениях итог работы.

Приложения с лекции 2 нет, так как я сделал скриншот в буфер обмена и забыл его сохранить, приступив к правкам своего отчёта, поэтому скриншот был благополучно утерян. 
Думаю, что моё посещение лекции можно замерить посредстам проверке моей работы.

По окончании данной лекции были сформированы пункты, которые нужно было исправить в моём отчёте.

\subsection{Разработка шаблона для ВКР}

При разработке шаблона я опирался на документ «Требования к выпускным квалификационным работам» Университета ИТМО (ЛНАОБУЧ-СМК-03-05-2022, версия 4.0) 
и использовал требования к оформлению и структуре документа.
В частности, я учитывал требования к формату A4, полуторному интервалу, шрифту Times New Roman 14~пт, 
полям (левое 30~мм, правое 15~мм, верхнее и нижнее по 20~мм), абзацному отступу 1{,}25~см, 
нумерации страниц арабскими цифрами по центру внизу страницы, а также оформлению заголовков и правилу начала каждой главы с новой страницы.

Также при реализации я ориентировался на разделы документа: 4.3 «Нумерация страниц ВКР», 4.4 «Оформление заголовков», 4.5 «Оформление иллюстраций», 4.6 «Оформление таблиц», 4.7 «Оформление формул и уравнений», 4.8 «Оформление списка сокращений и условных обозначений», 4.9 «Оформление списка терминов», 4.10 «Оформление списка использованных источников», 4.11 «Оформление приложений» (со всеми подразделами).

Отдельно я учитывал конкретные подпункты требований и переносил их в настройки шаблона, чтобы оформление было единообразным и проверяемым:
\begin{itemize}
  \item (4.3.1) Я настроил сквозную нумерацию страниц и вывод номера по центру нижнего колонтитула тем же шрифтом, что и основной текст (через настройку колонтитулов).
  \item (4.4.1--4.4.4) Я зафиксировал оформление заголовков и структурных элементов: ненумеруемые элементы выводились по центру прописными буквами, а главы начинались с новой страницы; нумеруемые разделы и подразделы оформлялись единообразно средствами пакета для заголовков.
  \item (4.5.2--4.5.4) Для иллюстраций я настроил подписи и нумерацию внутри разделов, а также использование ссылок на рисунки в тексте, чтобы обеспечить корректные перекрёстные ссылки.
  \item (4.6.1--4.6.3) Для таблиц я зафиксировал формат подписи «Таблица N — Наименование», размещение подписи сверху и возможность переноса длинных таблиц на следующие страницы.
  \item (4.7.1--4.7.4) Для формул я настроил оформление через окружения формул с нумерацией по разделам и добавлением пояснений под формулой, чтобы соблюдалась требуемая структура записи.
  \item (4.11.1--4.11.3) Для приложений я реализовал оформление с началом каждого приложения с новой страницы, буквенным обозначением и сохранением сквозной нумерации страниц по всему документу.
\end{itemize}

Для реализации шаблона я выбрал LaTeX, так как он позволял централизованно зафиксировать требования к оформлению в виде кода и переиспользовать их для любых документов, 
а также гарантировать воспроизводимость верстки при изменении содержания.
Практически это было реализовано через подключение и настройку пакетов (например, fancyhdr для нумерации страниц, titlesec для заголовков, caption для подписей рисунков и таблиц, amsmath для формул, longtable/ltablex для длинных таблиц, hyperref для ссылок), а также через собственные команды и правила оформления, вынесенные в файл стилей.
Шаблон стилей я реализовал в виде отдельного файла стилей с подключением через основной документ.
Геометрию страницы я настроил через пакет geometry, форматирование заголовков --- через titlesec.
Для генерации содержания я использовал пакет tocloft с заполнением точками между названием и номером страницы.
Оформление иллюстраций и таблиц я реализовал через caption (подписи «Рисунок» и «Таблица», разделитель-тире, нумерация внутри разделов).
Поддержку длинных таблиц я реализовал через longtable и ltablex.
Оформление формул я выполнил через amsmath с нумерацией внутри разделов и отдельной командой для автоматического добавления пояснений с перечислением переменных.
Оформление списков я настроил через enumitem (маркированные списки с тире и отступы согласно требованиям).
Для приложений я реализовал команды \slash CUSTOMappendixsetup (переключение нумерации на кириллические буквы: А, Б, В, Г...) и \slash CUSTOMappendix (создание приложения с добавлением в содержание).
Библиографию я переопределил через окружение thebibliography.

В итоге я получил рабочий шаблон ВКР в LaTeX в виде набора правил оформления и команд, который позволял собирать документ с соблюдением требований к нумерации, заголовкам, рисункам, таблицам, формулам, спискам и приложениям. Трудностей при работе с LaTeX не было.


\subsection{Техническое задание по ВКР}

При подготовке технического задания я ориентировался на требования к структуре 
и оформлению из документа «Требования к выпускным квалификационным работам» Университета ИТМО (ЛНАОБУЧ-СМК-03-05-2022, версия 4.0), 
чтобы итоговый документ соответствовал единому стилю и правилам оформления.

Техническое задание я составил для проекта выпускной квалификационной работы с названием «Разработка гибридной системы интеллектуального нормоконтроля с интеграцией LLM для семантического анализа технической документации».
При составлении ТЗ я опирался на следующие источники: Кузнецова~Е.\,С., Федюнина~Е.\,А., Орлов~И.\,А., Чукин~П.\,Е. Разработка алгоритма процесса согласования документов в 1С: PLM // Электронные информационные системы. 2022. №~1 (32). С.~48--58. EDN: QWTYLC; Баданина~Н.\,Д., Судаков~В.\,А., Яшин~Н.\,А. Программный комплекс визуализации и моделирования на основе BPMN нотации // Научная визуализация. 2022. Т.~14. №~3. С.~13--28. DOI: 10.26583/sv.14.3.02. EDN: EQVPDK; Шкуропадский~И.\,В., Овчинников~П.\,В., Ткачев~А.\,Н. Разработка систем интеллектуальной автоматизации бизнес-процессов на основе технологий машинного обучения // Друкеровский вестник. 2021. №~6 (44). С.~45--56. DOI: 10.17213/2312-6469-2021-6-45-56.

При наполнении разделов я действовал так.
Для функциональных требований сначала начертил use case диаграмму и затем вывел из неё формулировки требований.
Назначение и актуальность проблем текущей ситуации я определял через опрос.
Структуру разрабатываемой системы я сначала нарисовал в Eraiser: расставил связи и потоки данных, после чего перенёс это в ТЗ в виде текста; описание структурных модулей делал, анализируя схему --- потоки данных и связи между компонентами.
Описание пользовательского интерфейса вывел из функциональных требований; нефункциональные требования (надёжность, производительность и др.) --- из архитектуры приложения по схеме и отталкиваясь от метрик, которые называли респонденты при опросе.
Стадии и этапы разработки я сформировал, перебрав нарисованную архитектуру и разделив поэтапно интеграцию компонентов и шаги по подготовке и защите ВКР.
В итоге я последовательно сформировал разделы ТЗ: наименование (краткое и полное), назначение, основание для разработки, функциональные требования, структура разрабатываемой системы, пользовательский интерфейс, требования к надёжности и безопасности, условия эксплуатации, документация, стадии и этапы разработки, порядок контроля и приёмка.

Для оформления ТЗ я также использовал LaTeX по тем же соображениям, что и при разработке шаблона ВКР: мне было важно зафиксировать оформление в виде кода и получить воспроизводимый результат при правках содержания.
Кроме того, я переиспользовал отдельные элементы шаблона (настройки шрифтов, полей, заголовков и нумерации), чтобы сохранить одинаковый стиль и структуру оформления документов.

Итоговый документ технического задания в формате PDF приведён в приложении~Б. Трудностей при работе с LaTeX не было.
